\documentclass[10pt,a4paper]{amsart}
\usepackage[margin=19mm]{geometry}
\usepackage{amsmath,amssymb,mathtools}
\usepackage[scr=rsfs,cal=euler]{mathalfa}
\usepackage{xfrac}
\usepackage[unicode,hidelinks,bookmarks=false]{hyperref}
\newcommand{\ie}{i.e.\ }
\newcommand{\cf}{cf.\ }
\newcommand{\resp}{respectively\ }
\newcommand{\Subsection}{\subsection}
\providecommand{\nobreakdash}{\nobreak\hbox{-}}
\setlength{\emergencystretch}{2em}
\multlinegap0pt
\usepackage{bm}
\usepackage{bbm}
\usepackage{dsfont}
\usepackage{xspace}
\usepackage{cancel}
\usepackage{mathtools}
\usepackage{amsthm}
\makeatletter
\@ifundefined{theorem}{\newtheorem{theorem}{Theorem}}{}
\@ifundefined{lemma}{\newtheorem{lemma}[theorem]{Lemma}}{}
\makeatother
\title[Quantum recurrences and the arithmetic of Floquet dynamics]{Quantum recurrences and the arithmetic of Floquet dynamics}
\author{Amit Anand, Dinesh Valluri, Jack Davis, Shohini Ghose}
\date{Source: arXiv:2508.09933v3 (CC BY 4.0)}
\begin{document}
\maketitle

\noindent\emph{Source and licence.} Quantum recurrences and the arithmetic of Floquet dynamics. Version arXiv:2508.09933v3 (CC BY 4.0), distributed under the Creative Commons Attribution 4.0 International licence, as recorded in the arXiv licence metadata for this exact version and shown on the abstract page https://arxiv.org/abs/2508.09933v3. Full rights notice: licenses/arxiv-2508.09933-CC-BY-4.0.txt.\par\smallskip

\noindent\textit{Editorial scope.} Counted unit: the Lemma whose source label is \texttt{lemma:lemma 1} in the article (source0.tex lines 195--201, source label \texttt{lemma:lemma 1}), delivered together with the article's own complete original proof of it (source0.tex lines 202--245, one \texttt{proof} environment of 44 lines). No mathematics is written, repaired or re-derived: statement and proof are the article's own text line for line. Required context retained uncounted: none. Every definition, display and label of those lines is retained unchanged. Omitted from this package and not redistributed: 1--194, 246--493. Nothing inside a retained excerpt is omitted or altered: every retained body is delivered exactly as the article's lines are written, so no omission receipt of any kind accompanies this package. Citations: the retained excerpts carry no citation command at all, so no bibliography entry of the article is transcribed and no reference is redistributed. Reference adaptation: none was needed and none was made; every reference inside the delivered unit resolves inside it, so no unresolved reference or citation is produced. Printed slips kept literal: no printed word, symbol, hypothesis or number of the counted statement or of its proof was altered. Layout and class: the article's own class is \texttt{quantumarticle} with packages including inputenc, fontenc, babel, amsmath, amssymb, amsfonts, amsthm, mathtools, braket, biblatex, xurl, hyperref, setspace, graphicx; the wrapper is the lane's pinned \texttt{amsart} editorial wrapper at 10pt on A4, which restates the theorem-like environments the retained text needs and adds the article's own macro definitions that the retained text needs (none), with extra wrapper packages (bm, bbm, dsfont, xspace, cancel, mathtools). The delivered numbering is the unit's own. Assets: no retained span contains a figure environment, a graphics-inclusion command, a PDF-page inclusion, a table or a TikZ picture; no image file is copied into this package, the package declares no asset dependency and no image file is redistributed with it. Licence: version arXiv:2508.09933v3 of this article is distributed under the Creative Commons Attribution 4.0 International licence, as recorded in the arXiv licence metadata for this exact version and shown on the abstract page https://arxiv.org/abs/2508.09933v3. A full rights notice is delivered at licenses/arxiv-2508.09933-CC-BY-4.0.txt. Characters: every retained excerpt is pure ASCII, so the delivered engine, XeLaTeX with its default Latin Modern text font, reports no missing character.
\begin{lemma} \label{lemma:lemma 1}
    Suppose $U$ is a $d\times d$ periodic matrix with period $n$ and $\lambda_{i} = \zeta_{n}^{k_i} \tau^{1/n}$ are the eigenvalues of $U$ as described above. Then the following are true: 
    \begin{itemize}
        \item[(a)] $R := gcd(k_1 - k_d, \ldots, k_{d-1} - k_d)$ is invertible modulo $n$. In other words, $gcd(R,n) = 1$.
        \item[(b)] There exist integers $a_1, \ldots a_d$ such that $a_1 + \ldots + a_d = 0$ and $a_1 k_1 + \ldots + a_d k_d = 1 \pmod{n}$. 
    \end{itemize}
\end{lemma}

\begin{proof}\label{proof:lemma 1}
     (a.) Suppose $R = gcd(k_1 - k_d, \ldots, k_{d-1} - k_d)$, in particular $k_i \equiv k_d \pmod{R}$ for all $1 \leq i \leq d-1.$  write $k_i = k_{d} + Rl_{i}$ for some integers $l_i,$ for all $1 \leq i \leq d-1.$ We need to prove that $gcd(R,n) =1.$
     Observe,
     $\zeta_{n}^{k_i} = \zeta_{n}^{k_d} \zeta_{n}^{Rl_i} $. Since $\lambda_{i} = \zeta_{n}^{k_i} \tau^{1/n}$, we have 
     \begin{align}\label{eq: ratio of two eigne value}
         \frac{\lambda_i}{\lambda_d} &= \frac{\zeta_{n}^{k_i}}{ \zeta_{n}^{k_d}} = \zeta_{n}^{Rl_i}.
     \end{align}
     By raising both sides of Eq. \ref{eq: ratio of two eigne value} to the power of $n/gcd(R,n)$ we get
     \begin{align}
          \bigg(\frac{\lambda_i}{\lambda_d}\bigg)^{n/gcd(R,n)} = \bigg(\zeta_{n}^n\bigg)^{l_i\frac{R}{gcd(R,n)}} = 1,
     \end{align}
    since $\zeta_{n}$ is a $n$-th root of unity and  $Rl_i/gcd(R,n)$ is an integer. It follows that
\begin{equation}\label{eq:diag-power}
\begin{pmatrix}
\lambda_{1} &        &        &        \\
            & \ddots &        &        \\
            &        & \ddots &        \\
            &        &        & \lambda_{d}
\end{pmatrix}^{\tfrac{n}{\gcd(R,n)}}
\;=\;
\begin{pmatrix}
\lambda_{d} &        &        &        \\
            & \ddots &        &        \\
            &        & \ddots &        \\
            &        &        & \lambda_{d}
\end{pmatrix}^{\tfrac{n}{\gcd(R,n)}}
\;=\;
\lambda_{d}^{\tfrac{n}{\gcd(R,n)}}\,I.
\end{equation}
This implies $U^{n/gcd(R,n)} = \tau' I $ for $\tau' = \lambda_d^{n/gcd(R,n)} \in \mathrm{U}(1)$. This is only possible when $gcd(R,n) = 1,$ as $U$ has period $n$.

(b.) By Bezout's identity there exists integers $a_1, \ldots, a_{d-1}$ such that 
\begin{equation}\label{eq: writing Bezout identity for R}
    R := gcd(k_1 - k_d, \ldots, k_{d-1} - k_d) = b_1(k_1 - k_{d}) + \ldots + b_{d-1} (k_{d-1} - k_d) 
\end{equation}
for some integer $b_i$'s.  Define $b_d := -(b_1 +\ldots + b_{d-1}),$ and let $S$ be the inverse of $R$ modulo $n$, i.e., $RS = 1 \pmod{n}$. By multiplying $S$ on both sides of Eq. \eqref{eq: writing Bezout identity for R}, we get $Sb_1(k_1 - k_{d}) + \ldots + Sb_{d-1} (k_{d-1} - k_d) = RS = 1 \pmod{n}$. By defining $a_i :=Sb_i$ for $1\leq i \leq d$, we get 
\begin{align}
    a_1(k_1 - k_{d}) + \ldots + a_{d-1} (k_{d-1} - k_d) &= 1 \pmod{n},  \, \,  \text{since}\\
     \nonumber a_1(k_1 - k_{d}) + \ldots + a_{d-1} (k_{d-1} - k_d) & = a_1k_1 + \ldots + a_{d-1}k_{d-1} - ( a_1 + \dots + a_{d-1})k_d \\
   \nonumber  & = a_1k_1 + \ldots + a_dk_d = 1 \pmod{n},
\end{align}

as required in the lemma. 
\end{proof}

\end{document}
